\section{The auxiliary powers in the theta-rule DVR test}
\label{sec:key-rule-dvr}

The good-reduction step in Section~\ref{sec:key-different} uses a
unit normalization for the theta scalar on an integral polarized
abelian family. We explain how the local construction can give
unit values at \emph{every} residue characteristic even though
the construction of an individual powered invariant uses an
auxiliary integer which must be invertible. The word level has
two different roles here. The auxiliary integer in this argument
is chosen for a determinant calculation on sections of the
polarizing line. The fixed level of the moduli problem satisfies
the separate condition in Section~\ref{sec:key-different}.

\subsection{Choosing transverse-intersection degrees at a given prime}

For the numerical construction let \(d\geq1\) denote the dimension
parameter for the intersection. Given a prime \(p\), set
\[
       n=4p,\qquad b=dn+1,\qquad b'=b+p,\qquad
       P=\prod_{j=0}^d (b-jn)(b'-jn).
\]
Then \(n\ge4\) and \(dn<b<b'\). Each factor in \(P\) is \(1\bmod p\),
so \(P>0\) and \(p\nmid P\), including at \(p=2\). The exponent in
the determinant construction is
\[
     c=2P^{3d+2}\,d!\,n^d(b'-b)>0 .
\]
It can be divisible by \(p\). Only \(P\) is required to be a unit
in this part of the construction.

For every \(m\geq0\) choose a four-squares representation
\(m=a^2+b_1^2+b_2^2+b_3^2\). The integer matrix
\[
 U_m=
 \begin{pmatrix}
 a&-b_1&-b_2&-b_3\\
 b_1&a&-b_3&b_2\\
 b_2&b_3&a&-b_1\\
 b_3&-b_2&b_1&a
 \end{pmatrix}
\]
has \(U_m^{\,t}U_m=mI_4\): the columns have squared norm \(m\)
and their pairwise dot products cancel termwise. These matrices
supply the square-determinant comparisons at degrees \(b-jn\)
and \(b'-jn\). This proves all the arithmetic admissibility
conditions.

Here is the geometric input used with these integers. A tuple of
\(d\) sufficiently general sections of \(L^n\) on the special
fiber cuts out \(d!n^d\) transverse points for the principal
polarization. The release picks a basis of \(H^0(L^n)\), gives
the tuple algebraically independent coefficients over the residue
field, and passes to a local injection into a complete DVR in
which the finitely many coefficients have lifts. Its transversality
theorem gives the required étale zero locus on the extended residue
field. After another injective local extension, the determinant
comparison theorem supplies a unit from the two section Koszul
complexes, their evaluation at those points, the Jacobian product,
and the comparisons on fourth powers given by the \(U_m\). This
step requires \(P\) to be invertible. The separate theorems on
independence of the sections and change of base make the resulting
powered invariant functorial. These intersection, determinant,
and functoriality assertions are geometric inputs to the following
valuation argument.
\footnote{\href{https://github.com/ChristianKudeba/alpoge-rh-verification/blob/alpoge-source-7cbf7a0c/src/Development/AbelianVarietyTheory/Proofs/ThetaLawProofs/ExistsAdmissibleKeyNumericalDataPrimeTo.lean}{Admissible prime-to-\(p\) parameters and the quaternions};
\href{https://github.com/ChristianKudeba/alpoge-rh-verification/blob/alpoge-source-7cbf7a0c/src/Development/AbelianVarietyTheory/Definitions/KeyFormula/DegProd.lean}{degrees and powers in the determinant invariant};
\href{https://github.com/ChristianKudeba/alpoge-rh-verification/blob/alpoge-source-7cbf7a0c/src/Development/AbelianVarietyTheory/Proofs/ThetaLawProofs/ExistsIntegralKeyInvariantRuleCharZero.lean}{residue-field extension, transverse intersections and comparison}.}

\subsection{Changing the auxiliary integer}

In abelian dimension \(g\ge1\), the normalization step uses the
previous numerical construction at \(d=2g\). For each prime \(p\)
it produces a positive exponent \(N_p=2c\), an integer \(P_p\)
prime to \(p\), and a scalar \(F_p(x)\) functorial on framed
data \(x\) over fields of characteristic zero. At a DVR where
\(P_p\) is a unit its value on the generic fiber of integral
data is a unit. On a complex datum in normalized half-characteristic
coordinates it has value
\[
          F_p(x)=\bigl(\det(M)^4/c_0^8\bigr)^{N_p},
\]
where \(c_0\) is the nonzero theta scaling. These assertions are
the output of the normalized determinant construction just invoked.

For any two auxiliary primes \(p,q\), and any field \(K\) of
characteristic zero, one has
\[
                         F_p(x)^{N_q}=F_q(x)^{N_p}.
\]
To see this, descend the framed datum to a countable subfield
\(K_0\subset K\), as permitted by finite data descent in the
source. Embed \(K_0\) into \(\mathbb C\) and extend the datum
there. The two sides are equal to the same power of
\(\det(M)^4/c_0^8\), by the complex uniformization. Functoriality
and injectivity of \(K_0\hookrightarrow\mathbb C\) give equality
over \(K_0\). Extending scalars back gives the formula over \(K\).

Fix one choice \(F_2\). In residue characteristic zero \(P_2\)
is invertible. If the residue characteristic is \(p>0\) and
\(P_2\) is a nonunit, use \(F_p\), for which \(P_p\) is a unit.
For a DVR valuation \(v\) on \(K\) and the generic fiber of the
integral frame,
\[
           N_p v(F_2(x))=N_2 v(F_p(x))=0 .
\]
Since \(N_p>0\), this gives \(v(F_2(x))=0\).
Thus the fixed powered rule is a unit at all residue
characteristics. This computation uses no \(p\)-adic division by
the integer \(N_p\) in the valuation ring.

The exponent-one scalar used in the manuscript also requires
a coherent root over arbitrary characteristic-zero fields.
The complex candidate \(t(x)=\det(M)^4/c_0^8\) is fixed by the
theta transformation law. The release establishes its descent
using connected families of polarized varieties, equivariance of
the complex value, and its specialization on smooth curves. The
elementary uniqueness step is as follows. Suppose a smooth
integral complex base carries a unit \(u\) and a continuous
normalized value \(t\) with \(t^N=u^N\) on the complex points.
Their ratio takes values in the finite discrete group \(\mu_N\);
the complex points of this base are connected, so the ratio is
constant. An equality \(u=t\) in one fiber identifies them on
the entire base. Establishing the indicated families and their
arithmetic comparison is additional to this observation.
Once descent gives the exact rule \(\kappa\) functorial in the
field and with the displayed normal value, apply the same
countable-field argument to obtain
\(\kappa(x)^{N_2}=F_2(x)\). Taking valuations gives
\(v(\kappa(x))=0\), proving the unit conclusion for an integral
proper theta datum at every residue characteristic under the
stated inputs.
\footnote{\href{https://github.com/ChristianKudeba/alpoge-rh-verification/blob/alpoge-source-7cbf7a0c/src/Development/AbelianVarietyTheory/Proofs/ThetaLawProofs/ThetaKeyRuleExistsOfAlgebraicInvariants.lean}{Cross-power equality and replacement of the auxiliary integer};
\href{https://github.com/ChristianKudeba/alpoge-rh-verification/blob/alpoge-source-7cbf7a0c/src/Development/AbelianVarietyTheory/Proofs/ThetaLawProofs/ThetaFrameValueEqUnitSmoothConnectedBase.lean}{constancy of the root on smooth connected bases};
\href{https://github.com/ChristianKudeba/alpoge-rh-verification/blob/alpoge-source-7cbf7a0c/src/Development/AbelianVarietyTheory/Proofs/UniformizationProofs/KeyFormulaOfConnFam.lean}{the geometric and specialization inputs for the exponent-one rule}.}
